% --- Archon physics formalization source begin ---
% archon:physics
% archon:covers PhyXMiniProblems/problem_phyx_mini_0342.lean
% archon:source-report reports/phyx_mini/problem_phyx_mini_0342.source.json

\chapter{Physics problem phyx\_mini\_0342}
\label{ch:PhyXMiniProblems_problem_phyx_mini_0342}

\paragraph{Problem source.}
\#\# Physical scenario

Three moles of an ideal gas are taken around cycle \$acb\$ shown in \textbackslash{}textbf\{figure\}. For this gas, \$C\_p = 29.1\textasciitilde{}\textbackslash{}text\{J/mol\} \textbackslash{}cdot \textbackslash{}text\{K\}\$. Process \$ac\$ is at constant pressure, process \$ba\$ is at constant volume, and process \$cb\$ is adiabatic. The temperatures of the gas in states \$a\$, \$c\$, and \$b\$ are \$T\_a = 300\textasciitilde{}\textbackslash{}text\{K\}\$, \$T\_c = 492\textasciitilde{}\textbackslash{}text\{K\}\$, and \$T\_b = 600\textasciitilde{}\textbackslash{}text\{K\}\$.

\#\# Question

Calculate the total work \$W\$ for the cycle.

\#\# Answer choices

- A: A: -1.65 \textbackslash{}times 10\textasciicircum{}\{3\}J

- B: B: -1.95 \textbackslash{}times 10\textasciicircum{}\{3\}J

- C: C: 1.95 \textbackslash{}times 10\textasciicircum{}\{3\}J

- D: D: -1.45 \textbackslash{}times 10\textasciicircum{}\{3\}J

\#\# Figure caption (auxiliary; use the image as primary evidence)

The image is a pressure-volume (\textbackslash{}(p\textbackslash{})-\textbackslash{}(v\textbackslash{})) graph featuring three points: \textbackslash{}(a\textbackslash{}), \textbackslash{}(b\textbackslash{}), and \textbackslash{}(c\textbackslash{}). 

- The axes are labeled \textbackslash{}(p\textbackslash{}) (vertical) and \textbackslash{}(v\textbackslash{}) (horizontal), with their intersection marked as \textbackslash{}(O\textbackslash{}).
- Three black points labeled \textbackslash{}(a\textbackslash{}), \textbackslash{}(b\textbackslash{}), and \textbackslash{}(c\textbackslash{}) are connected by blue arrows, forming a closed loop.
  - The arrow from \textbackslash{}(a\textbackslash{}) to \textbackslash{}(b\textbackslash{}) is vertical.
  - The arrow from \textbackslash{}(b\textbackslash{}) to \textbackslash{}(c\textbackslash{}) is curved.
  - The arrow from \textbackslash{}(c\textbackslash{}) to \textbackslash{}(a\textbackslash{}) is horizontal.
- The blue arrows indicate the direction of the cycle from \textbackslash{}(a\textbackslash{}) to \textbackslash{}(b\textbackslash{}) to \textbackslash{}(c\textbackslash{}) and back to \textbackslash{}(a\textbackslash{}).

This likely represents a cyclic process in a thermodynamic system.

\#\# Dataset metadata

Laws of Thermodynamics; Multi-Formula Reasoning, Physical Model Grounding Reasoning

\paragraph{Recorded answer/context.}
B: B: -1.95 \textbackslash{}times 10\textasciicircum{}\{3\}J

\paragraph{Figure/image path.}
/root/proposal\_for\_physic/science-mango/phyx\_mini\_run/phyx\_data/test\_image/342.png

\paragraph{Formalization target.}
create a compiling Lean file with sorry bodies at `PhyXMiniProblems/problem\_phyx\_mini\_0342.lean`. The Lean declarations must preserve the physical quantities, dimensions or dimensional roles, figure labels, governing-law hypotheses, and final relation expressed by this problem.
Use Mathlib/Physlib names found through LeanExplore where available. If a physics API is missing, introduce faithful local abstractions rather than scalar placeholder aliases.

\begin{theorem}[Physics formalization target]
\label{thm:physics:phyx_mini_0342:target}
\lean{PhyXMiniProblems.ProblemPhyXMini0342.totalWorkDoneByGas_rounds_to_negative_1Point95_kilojoules}
\uses{def:physics:phyx-mini-0342:phyxminiproblems-problemphyxmini0342-idealgascyclesetup, def:physics:phyx-mini-0342:phyxminiproblems-problemphyxmini0342-hasphysicalidealgascycleparameters, def:physics:phyx-mini-0342:phyxminiproblems-problemphyxmini0342-satisfiesidealgascyclelaws, def:physics:phyx-mini-0342:phyxminiproblems-problemphyxmini0342-hasgivenidealgascycledata, def:physics:phyx-mini-0342:phyxminiproblems-problemphyxmini0342-roundstonearesttenjoules}
The assigned autoformalize agent should translate this physics problem into Lean declarations in the covered file, with theorem and lemma proof bodies written as `by sorry`.
\end{theorem}
\begin{proof}
This is an autoformalization task, not a proof task. Produce statements that can later be proved by the physics prover without weakening the physical model.
\end{proof}

% --- Archon declaration topology begin ---
\section{Declaration topology}
\begin{definition}[Gas Volume]
  \label{def:physics:phyx-mini-0342:phyxminiproblems-problemphyxmini0342-gasvolume}
  \lean{PhyXMiniProblems.ProblemPhyXMini0342.GasVolume}
  A physical gas volume, with dimension L³
\end{definition}
\begin{proof}
  This is the declared model definition of Gas Volume.
\end{proof}
\begin{definition}[pressure In Pascals]
  \label{def:physics:phyx-mini-0342:phyxminiproblems-problemphyxmini0342-pressureinpascals}
  \lean{PhyXMiniProblems.ProblemPhyXMini0342.pressureInPascals}
  Pascal readout of a physical pressure
\end{definition}
\begin{proof}
  This is the declared model definition of pressure In Pascals.
\end{proof}
\begin{definition}[volume In Cubic Meters]
  \label{def:physics:phyx-mini-0342:phyxminiproblems-problemphyxmini0342-volumeincubicmeters}
  \lean{PhyXMiniProblems.ProblemPhyXMini0342.volumeInCubicMeters}
  \uses{def:physics:phyx-mini-0342:phyxminiproblems-problemphyxmini0342-gasvolume}
  Cubic-metre readout of a physical volume
\end{definition}
\begin{proof}
  This is the declared model definition of volume In Cubic Meters.
\end{proof}
\begin{definition}[energy In Joules]
  \label{def:physics:phyx-mini-0342:phyxminiproblems-problemphyxmini0342-energyinjoules}
  \lean{PhyXMiniProblems.ProblemPhyXMini0342.energyInJoules}
  Joule readout of a physical energy
\end{definition}
\begin{proof}
  This is the declared model definition of energy In Joules.
\end{proof}
\begin{definition}[Cycle State]
  \label{def:physics:phyx-mini-0342:phyxminiproblems-problemphyxmini0342-cyclestate}
  \lean{PhyXMiniProblems.ProblemPhyXMini0342.CycleState}
  The three labeled thermodynamic states in the supplied figure
\end{definition}
\begin{proof}
  This is the declared model definition of Cycle State.
\end{proof}
\begin{definition}[Cycle Leg]
  \label{def:physics:phyx-mini-0342:phyxminiproblems-problemphyxmini0342-cycleleg}
  \lean{PhyXMiniProblems.ProblemPhyXMini0342.CycleLeg}
  The directed legs, in the order indicated by the arrows in the image
\end{definition}
\begin{proof}
  This is the declared model definition of Cycle Leg.
\end{proof}
\begin{definition}[initial State]
  \label{def:physics:phyx-mini-0342:phyxminiproblems-problemphyxmini0342-cycleleg-initialstate}
  \lean{PhyXMiniProblems.ProblemPhyXMini0342.CycleLeg.initialState}
  \uses{def:physics:phyx-mini-0342:phyxminiproblems-problemphyxmini0342-cyclestate, def:physics:phyx-mini-0342:phyxminiproblems-problemphyxmini0342-cycleleg}
  Initial state of each directed leg
\end{definition}
\begin{proof}
  This is the declared model definition of initial State.
\end{proof}
\begin{definition}[final State]
  \label{def:physics:phyx-mini-0342:phyxminiproblems-problemphyxmini0342-cycleleg-finalstate}
  \lean{PhyXMiniProblems.ProblemPhyXMini0342.CycleLeg.finalState}
  \uses{def:physics:phyx-mini-0342:phyxminiproblems-problemphyxmini0342-cyclestate, def:physics:phyx-mini-0342:phyxminiproblems-problemphyxmini0342-cycleleg}
  Final state of each directed leg
\end{definition}
\begin{proof}
  This is the declared model definition of final State.
\end{proof}
\begin{definition}[Axis Quantity]
  \label{def:physics:phyx-mini-0342:phyxminiproblems-problemphyxmini0342-axisquantity}
  \lean{PhyXMiniProblems.ProblemPhyXMini0342.AxisQuantity}
  Physical quantities that can label an axis of the supplied diagram
\end{definition}
\begin{proof}
  This is the declared model definition of Axis Quantity.
\end{proof}
\begin{definition}[Figure Point Label]
  \label{def:physics:phyx-mini-0342:phyxminiproblems-problemphyxmini0342-figurepointlabel}
  \lean{PhyXMiniProblems.ProblemPhyXMini0342.FigurePointLabel}
  Labels visibly attached to the origin and the three plotted states
\end{definition}
\begin{proof}
  This is the declared model definition of Figure Point Label.
\end{proof}
\begin{definition}[Process Kind]
  \label{def:physics:phyx-mini-0342:phyxminiproblems-problemphyxmini0342-processkind}
  \lean{PhyXMiniProblems.ProblemPhyXMini0342.ProcessKind}
  Thermodynamic classification of a process leg
\end{definition}
\begin{proof}
  This is the declared model definition of Process Kind.
\end{proof}
\begin{definition}[Leg Shape]
  \label{def:physics:phyx-mini-0342:phyxminiproblems-problemphyxmini0342-legshape}
  \lean{PhyXMiniProblems.ProblemPhyXMini0342.LegShape}
  Geometric appearance of a directed process in the pressure-volume plot
\end{definition}
\begin{proof}
  This is the declared model definition of Leg Shape.
\end{proof}
\begin{definition}[Cycle Direction]
  \label{def:physics:phyx-mini-0342:phyxminiproblems-problemphyxmini0342-cycledirection}
  \lean{PhyXMiniProblems.ProblemPhyXMini0342.CycleDirection}
  Orientation of the closed loop read from the arrowheads in the image
\end{definition}
\begin{proof}
  This is the declared model definition of Cycle Direction.
\end{proof}
\begin{definition}[Pressure Volume Figure]
  \label{def:physics:phyx-mini-0342:phyxminiproblems-problemphyxmini0342-pressurevolumefigure}
  \lean{PhyXMiniProblems.ProblemPhyXMini0342.PressureVolumeFigure}
  \uses{def:physics:phyx-mini-0342:phyxminiproblems-problemphyxmini0342-cyclestate, def:physics:phyx-mini-0342:phyxminiproblems-problemphyxmini0342-cycleleg, def:physics:phyx-mini-0342:phyxminiproblems-problemphyxmini0342-axisquantity, def:physics:phyx-mini-0342:phyxminiproblems-problemphyxmini0342-figurepointlabel, def:physics:phyx-mini-0342:phyxminiproblems-problemphyxmini0342-processkind, def:physics:phyx-mini-0342:phyxminiproblems-problemphyxmini0342-legshape, def:physics:phyx-mini-0342:phyxminiproblems-problemphyxmini0342-cycledirection}
  The labeled pressure-volume diagram accompanying the exercise
\end{definition}
\begin{proof}
  This is the declared model definition of Pressure Volume Figure.
\end{proof}
\begin{definition}[Thermodynamic State]
  \label{def:physics:phyx-mini-0342:phyxminiproblems-problemphyxmini0342-thermodynamicstate}
  \lean{PhyXMiniProblems.ProblemPhyXMini0342.ThermodynamicState}
  \uses{def:physics:phyx-mini-0342:phyxminiproblems-problemphyxmini0342-gasvolume}
  Pressure, volume, and absolute temperature at one equilibrium state
\end{definition}
\begin{proof}
  This is the declared model definition of Thermodynamic State.
\end{proof}
\begin{definition}[Ideal Gas Cycle Setup]
  \label{def:physics:phyx-mini-0342:phyxminiproblems-problemphyxmini0342-idealgascyclesetup}
  \lean{PhyXMiniProblems.ProblemPhyXMini0342.IdealGasCycleSetup}
  \uses{def:physics:phyx-mini-0342:phyxminiproblems-problemphyxmini0342-cyclestate, def:physics:phyx-mini-0342:phyxminiproblems-problemphyxmini0342-cycleleg, def:physics:phyx-mini-0342:phyxminiproblems-problemphyxmini0342-pressurevolumefigure, def:physics:phyx-mini-0342:phyxminiproblems-problemphyxmini0342-thermodynamicstate}
  The ideal gas, its three states, and the energetic quantities assigned to the three directed legs. The heat sign convention is heat transferred to the gas; the work sign convention is work done by the gas. amountOfGasMoles, the two molar heat capacities, and the gas constant are explicit scalar readouts in the units stated in their field names
\end{definition}
\begin{proof}
  This is the declared model definition of Ideal Gas Cycle Setup.
\end{proof}
\begin{definition}[Has Physical Ideal Gas Cycle Parameters]
  \label{def:physics:phyx-mini-0342:phyxminiproblems-problemphyxmini0342-hasphysicalidealgascycleparameters}
  \lean{PhyXMiniProblems.ProblemPhyXMini0342.HasPhysicalIdealGasCycleParameters}
  \uses{def:physics:phyx-mini-0342:phyxminiproblems-problemphyxmini0342-pressureinpascals, def:physics:phyx-mini-0342:phyxminiproblems-problemphyxmini0342-volumeincubicmeters, def:physics:phyx-mini-0342:phyxminiproblems-problemphyxmini0342-idealgascyclesetup}
  Positivity assumptions selecting physical ideal-gas states and constants
\end{definition}
\begin{proof}
  This is the declared model definition of Has Physical Ideal Gas Cycle Parameters.
\end{proof}
\begin{definition}[Satisfies Ideal Gas Cycle Laws]
  \label{def:physics:phyx-mini-0342:phyxminiproblems-problemphyxmini0342-satisfiesidealgascyclelaws}
  \lean{PhyXMiniProblems.ProblemPhyXMini0342.SatisfiesIdealGasCycleLaws}
  \uses{def:physics:phyx-mini-0342:phyxminiproblems-problemphyxmini0342-pressureinpascals, def:physics:phyx-mini-0342:phyxminiproblems-problemphyxmini0342-volumeincubicmeters, def:physics:phyx-mini-0342:phyxminiproblems-problemphyxmini0342-energyinjoules, def:physics:phyx-mini-0342:phyxminiproblems-problemphyxmini0342-idealgascyclesetup}
  Ideal-gas mechanics and the first law for the cycle, expressed in coherent SI readouts. The convention Q = ΔU + W uses W > 0 for work done by the gas
\end{definition}
\begin{proof}
  This is the declared model definition of Satisfies Ideal Gas Cycle Laws.
\end{proof}
\begin{definition}[Has Given Ideal Gas Cycle Data]
  \label{def:physics:phyx-mini-0342:phyxminiproblems-problemphyxmini0342-hasgivenidealgascycledata}
  \lean{PhyXMiniProblems.ProblemPhyXMini0342.HasGivenIdealGasCycleData}
  \uses{def:physics:phyx-mini-0342:phyxminiproblems-problemphyxmini0342-idealgascyclesetup}
  The numerical givens and trustworthy readouts from the primary image. In particular, the arrowheads give a → c → b → a, while horizontal and vertical alignment give equal pressure on a,c and equal volume on b,a
\end{definition}
\begin{proof}
  This is the declared model definition of Has Given Ideal Gas Cycle Data.
\end{proof}
\begin{definition}[Rounds To Nearest Ten Joules]
  \label{def:physics:phyx-mini-0342:phyxminiproblems-problemphyxmini0342-roundstonearesttenjoules}
  \lean{PhyXMiniProblems.ProblemPhyXMini0342.RoundsToNearestTenJoules}
  \uses{def:physics:phyx-mini-0342:phyxminiproblems-problemphyxmini0342-energyinjoules}
  The joule readout of work rounds to reportedJoules at the nearest-ten-joule precision implicit in a three-significant-figure value near two kilojoules
\end{definition}
\begin{proof}
  This is the declared model definition of Rounds To Nearest Ten Joules.
\end{proof}
% --- Archon declaration topology end ---
% --- Archon physics formalization source end ---
